\begin{thebibliography}{99}
\bibitem{source1} R. Blute, T. Ehrhard, and C. Tasson. A convenient differential category. Cahiers de topologie et géométrie différentielle catégoriques, 2012.
\bibitem{source2} R. F. Blute, J. R. B. Cockett, J.-S. P. Lemay, and R. A. G. Seely. Differential categories revisited. Applied Categorical Structures, 2020.
\bibitem{source3} R. F. Blute, J. R. B. Cockett, and R. A. G. Seely. Differential categories. Math. Struct. Comput. Sci., 2006.
\bibitem{source4} F. Breuvart, M. Kerjean, and S. Mirwasser. Unifying graded linear logic and differential operators. In 8th International Conference on Formal Structures for Computation and Deduction, FSCD, 2023.
\bibitem{source5} J. R. B. Cockett and R. A. G. Seely. Weakly distributive categories. Journal of Pure and Applied Algebra, 114(2), 1997.
\bibitem{source6} J. R. B. Cockett and R.A.G. Seely. Proof theory for full intuitionistic linear logic, bilinear logic, and mix categories. Theory and Applications of categories, 3(5), 1997.
\bibitem{source7} J. R. B. Cockett and R.A.G. Seely. Linearly distributive functors. Journal of Pure and Applied Algebra, 143(1-3), 1999.
\bibitem{source8} J.R.B. Cockett and P. V. Srinivasan. Exponential modalities and complementarity (extended abstract). In Kohei Kishida, editor, Proceedings of the Fourth International Conference on Applied Category Theory, Cambridge, United Kingdom, 12-16th July 2021, volume 372 of Electronic Proceedings in Theoretical Computer Science. Open Publishing Association, 2022.
\bibitem{source9} T. Ehrhard. On Köthe Sequence Spaces and Linear Logic. Mathematical Structures in Computer Science, 12(5), 2002.
\bibitem{source10} T. Ehrhard. Finiteness spaces. Mathematical Structures in Computer Science, 15(4), 2005.
\bibitem{source11} T. Ehrhard. An introduction to differential linear logic: proof-nets, models and antiderivatives. Mathematical Structures in Computer Science, 2017.
\bibitem{source12} T. Ehrhard and L. Regnier. Differential interaction nets. Theoretical Computer Science, 364(2), 2006.
\bibitem{source13} A. Fleury and C. Retoré. The mix rule. Mathematical Structures in Computer Science, 4(2), 1994.
\bibitem{source14} R. Gannoun, R. Hachaichi, H. Ouerdiane, and A. Rezgui. Un théorème de dualité entre espaces de fonctions holomorphes à croissance exponentielle. Journal of Functional Analysis, 171(1), 2000.
\bibitem{source15} J.-Y. Girard. Linear Logic. Theoretical Computer Science, 50(1), 1987.
\bibitem{source16} A. Grothendieck. Produits tensoriels topologiques et espaces nucléaires. Memoirs of the AMS, 16, 1966. Publisher: American Mathematical Society.
\bibitem{source17} M. Kerjean. A Logical Account for Linear Partial Differential Equations. In Proceedings of the 33rd Annual ACM/IEEE Symposium on Logic in Computer Science, LICS 2018, Oxford, UK, July 09-12, 2018, 2018.
\bibitem{source18} M. Kerjean and J.-S. P. Lemay. Taylor Expansion as a Monad in Models of DiLL. In 2023 38th Annual ACM/IEEE Symposium on Logic in Computer Science (LICS), 2023.
\bibitem{source19} O. Laurent. Etude de la polarisation en logique. Thèse de Doctorat, Université Aix-Marseille II, March 2002.
\bibitem{source20} J.-S. P. Lemay. Exponential functions in cartesian differential categories. Applied Categorical Structures, 29, 2021.
\bibitem{source21} J.-S. P. Lemay and J.-B. Vienney. Graded differential categories and graded differential linear logic. In Mathematical Foundations of Programming Semantics, 2023.
\bibitem{source22} P.-A. Melliès. Categorical semantics of linear logic. Société Mathématique de France, 2008.
\bibitem{source23} C.-H. L. Ong. Quantitative semantics of the lambda calculus: Some generalisations of the relational model. In 31st Annual ACM/IEEE Symposium on Logic in Computer Science (LICS), 2017.
\end{thebibliography}
